\documentclass[11pt]{article}
\usepackage{fontspec}
\usepackage[utf8]{inputenc}
\usepackage{xunicode}
\setmainfont{Helvetica}
\usepackage[paperwidth=11in,paperheight=8.5in,margin=1in,headheight=0.0in,footskip=0.5in,includehead,includefoot,portrait]{geometry}
\usepackage[absolute]{textpos}
\TPGrid[0.5in, 0.25in]{23}{24}
\parindent=0pt
\parskip=12pt
\usepackage{nopageno}
\usepackage{graphicx}
\graphicspath{ {./images/} }
\usepackage{amsmath}
\usepackage{tikz}
\usepackage{multicol}
\newcommand*\circled[1]{\tikz[baseline=(char.base)]{
            \node[shape=circle,draw,inner sep=1pt] (char) {#1};}}

\begin{document}

\begin{center}
\huge \pmb{Techniques}
\end{center}

\begin{multicols}{2}

\leftskip0.25in
\includegraphics[height=2cm]{./images/barriolage.png} \pmb{Quasi-barriolage} : This notation represents a static arrangement of fingers across the strings which is then slid up and down as notated by the contour of the glissando, simplified to show only the lowest note to reduce visual noise. The lower auxiliary staff shows strings crossing bowings which will produce a kind of irregular barriolage effect. \\
\rightskip\leftskip
\phantom{text} \hfill \phantom{()}

\leftskip0.25in
\includegraphics[height=1cm]{./images/behind_bridge.png} \pmb{Behind the bridge} : Note to be played behind the bridge are written with the relevant string tone and the abbreviation \textit{b.b.} where all such notes should be played on the wrapping of the string to produce a gritty, scratch sound. \\
\rightskip\leftskip
\phantom{text} \hfill \phantom{()}

\vspace*{0.25cm}

\leftskip0.25in
\includegraphics[height=1cm]{./images/circle_bow.png} \pmb{Circular bowing} : Circular bows should be performed at a speed where one full circle is completed over the note's duration which the symbol is over. \\
\rightskip\leftskip
\phantom{text} \hfill \phantom{()}

\vspace*{0.25cm}

\leftskip0.25in
\includegraphics[height=1cm]{./images/feather_beam.png} \pmb{Feathered beams} : Some accelerandi and ritardandi are notated by a feathered beam, widening the the direction of greater speed and narrowing in the direction of lesser speed. The full duration over which these gestures occur is notated in the surrounding bracket.\\
\rightskip\leftskip
\phantom{text} \hfill \phantom{()}

\vspace*{0.25cm}

\leftskip0.25in
\includegraphics[height=1cm]{./images/finger_pressure_transition.png} \pmb{Shifting finger pressure} : Continuous adjustments in the left hand finger pressure is sometimes notated in a spanner above the staff, but at other times is shown within the staff by an arrow attached to a glissando. In these instances, both the pitch should slide as in a normal glissando and simultaneously the finger pressure must change. \\
\rightskip\leftskip
\phantom{text} \hfill \phantom{()}

\vspace*{0.25cm}

\leftskip0.25in
\includegraphics[height=1cm]{./images/flared_hairpin.png} \pmb{Flared hairpin} : Some crescendi and diminuendi are notated with flared hairpins. These refer to a very subtle change in dynamic near the thin portion of the hairpin and a sudden change in dynamic near the flared portion of the hairpin. All other hairpins behave as typical crescendi and decrescendi. \\
\rightskip\leftskip
\phantom{text} \hfill \phantom{()}

\vspace*{0.25cm}

\leftskip0.25in
\includegraphics[height=1cm]{./images/highest_note.png} \pmb{Highest note} : An aleatoric high pitch, close to the highest pitch you can attain with the instrument, even including noise artifacts from such high notes, are written as an upward triangle. These notes need not be pure tones, however this symbol does no implicitly ask for scratch tone. \\
\rightskip\leftskip
\phantom{text} \hfill \phantom{()}

\vspace*{0.25cm}

\leftskip0.25in
\includegraphics[height=1cm]{./images/nail_flick.png} \pmb{Nail flick} :This symbol represents a percussive technique where the string is struck by the fingernail from a flicking motion. \\
\rightskip\leftskip
\phantom{text} \hfill \phantom{()}

\vspace*{0.25cm}

\leftskip0.25in
\includegraphics[height=1cm]{./images/nail_pluck.png} \pmb{Nail pluck} :This symbol represents an alternative pizzicato where the string is plucked with the fingernail like a plectrum. If the performer's nails are too short to comfortably achieve this technique, a typical pizzicato is acceptable. \\
\rightskip\leftskip
\phantom{text} \hfill \phantom{()}

\vspace*{0.25cm}

\leftskip0.25in
\includegraphics[height=1cm]{./images/pitch_bend.png} \pmb{Pitch bend} : A trailing glissando with no explicit termination tone should end approximately where the line ends relative to the staff however, precision is not important. \\
\rightskip\leftskip
\phantom{text} \hfill \phantom{()}

\vspace*{0.25cm}

\leftskip0.25in
\includegraphics[height=1.25cm]{./images/polyphony.png} \pmb{Polyphonic playing} : Moments of single-instrument polyphony are notated on a single staff. Prolation brackets apply only to the stems which face them. Upward stems are governed by brackets above and downward stems are governed by brackets below. The gestures within this piece approximately keep a constant central tone while the bow should alternate between the string above and the string below this note. [for a similar technique, see \textit{Tetras} by Iannis Xenakis at approximately 8:30] \\
\rightskip\leftskip
\phantom{text} \hfill \phantom{()}

\vspace*{0.25cm}

\leftskip0.25in
\includegraphics[height=1.25cm]{./images/scp_staff.png} \pmb{String contact point} : Bow position on the string is sometimes notated by an auxiliary staff where the higher, solid lines are degrees of \textit{sul ponticello} and the lower, dashed lines are degrees of \textit{sul tasto}. The thick line in the middle of the staff is a normal playing position. Whether these degrees of change are treated linearly or exponentially is left to the performers. Likewise, exactly where the most extreme tasto ends is left to the performers. In general, the composer imagines a linear treatment of the degrees with tasto ending about an inch across the finger board, but this is not a required reading. \\
\rightskip\leftskip
\phantom{text} \hfill \phantom{()}

\vspace*{0.25cm}

\leftskip0.25in
\includegraphics[height=1cm]{./images/shifting_tremolo.png} \pmb{Shifting tremolo} : Typically, all tremoli should be performed as fast as possible with smoothed bow changes. Tremoli should not be staccato. On occasion the speed of the tremolo should shift gradually as notated by tremolo slashes written above the staff. In this spanner, more lines refers to faster tremolo and fewer lines refers to slower tremolo. When tremolo slashes are present on the stem, the number of slashes is irrelevant. Sometimes fewer slashes are written on the stem in order to save space below a note with multiple beams. \\
\rightskip\leftskip
\phantom{text} \hfill \phantom{()}

\vspace*{0.25cm}

\leftskip0.25in
\includegraphics[height=1.25cm]{./images/staccato.png} \pmb{Staccato glissando} : This notation refers to a constant glissando with the bow playing staccato on the written rhythms. [see \textit{Sahara} by Francisco Guerrero Marín at approximately 1:20] \\
\rightskip\leftskip
\phantom{text} \hfill \phantom{()}

\vspace*{0.25cm}

\leftskip0.25in
\includegraphics[height=2cm]{./images/triple_trill.png} \pmb{Triple trill} : Sometimes more than one trill pitch is requested. For double trills, the only option is to trill to the first tone, return the main tone, then trill to the second tone, and repeat. Triple trills present a more complicated field of possible orders. Performers may improvise these trill figures in whatever order is preferred. \\
\rightskip\leftskip
\phantom{text} \hfill \phantom{()}

\vspace*{0.25cm}

\leftskip0.25in
\includegraphics[height=1cm]{./images/vibrato_contour.png} \pmb{Vibrato} : This piece should be performed primarily senza vibrato. Occasional vibrato to color (or ``warm'') the timbre of a tone is acceptable but should not be performed ``bravura.'' What an oscillating line is written above the staff, a move ``bravura'' style vibrato is requested and should follow the contour of the line. Narrow wavelengths represent a fast vibrato, wide wavelengths represent a slow vibrato, shallow wave amplitudes represent a narrow vibrato, and tall wave amplitudes represent a wide vibrato. \\
\rightskip\leftskip
\phantom{text} \hfill \phantom{()}

\vspace*{0.25cm}


\end{multicols}.

\end{document}
